\section{Stage 3: Explainable and Auditable Rankings}
\label{sec:lit-transparency}

The introduction of NLP-based screening has paved the way to a more accurate assessment of candidate skill in
many cases, however, this is at the cost of interpretability. Reviews of these algorithmic hiring systems describe
this as a ``transparency deficit''~\cite{bogen2018}. Often, recruiters are left with scores or ranks, with little
to no significant justification, significantly reducing trust in the system and also introducing
issues with fair-hiring compliance~\cite{raghavan2020,bogen2018}. The introduction of AI-assisted ATS systems
worsens this issue further since rankings are produced by opaque models trained on historical hiring data.

\subsection{Inspectability in LLM Pipelines}
Lo et al.~\cite{lo2025aihiring} aim to cross this gap when they separate the CV screening process into modular LLM stages.
They make the pipeline inspectable stage-by-stage, and report agreement with human HR ratings on anonymised resumes.
However, this modularity does not attribute the final score to specific evidence sources. Transparency in
hiring requires both understandable stages as well as attributable contributions from each source of evidence
to the outcome.

\subsection{Local Model-Agnostic Explanations}
The broader explainability literature offers techniques that are relevant when we design MERIT's glass-box layer.
Ribeiro et al.~\cite{ribeiro2016lime} introduce LIME (Local Interpretable Model-Agnostic Explanations), which approximates
a black-box prediction with a locally faithful linear model. In recruitment, this could show which CV keywords pushed a score
up or down. However, LIME attributions are local and can change across perturbation samples~\cite{ribeiro2016lime}. They also
explain \emph{features}, not \emph{sources of evidence}, so a recruiter might learn that ``Python'' mattered, but not whether that
came from a CV claim, GitHub commits, or a LinkedIn endorsement.

\subsection{Counterfactual Explanations}
Wachter et al.~\cite{wachter2017counterfactual} argue instead for counterfactual explanations, namely the closest
``what would have changed my outcome'' statement for the affected individual. In hiring, that might mean telling a rejected
candidate that higher GitHub activity in the past year would have changed the shortlist decision. Counterfactuals can be
actionable for candidates, however, they are hard to make plausible when evidence spans CV, GitHub, and LinkedIn, and some
facts such as past employment dates cannot be revised. They also note that GDPR Article~22 requires human intervention and
contestation rather than a legally binding duty to open the internal logic of the model~\cite{wachter2017counterfactual}.

\subsection{The Impact of Explanations on Fairness}
Dodge et al.~\cite{dodge2019explaining} highlight the importance of explanations, especially how the wording of an explanation is able
to shift an evaluator's fairness judgement, this concern is also repeated by Bogen and Rieke~\cite{bogen2018}, who 
argue that showing reasoning does not always produce fair outcomes, especially when the model or data that were employed were biased.
Rather than providing just explanations, MERIT provides full transparency, to allow recruiters to contest the information
when needed, rather than blindly trust an explanation.

\subsection{Motivation for Shapley Values}
The works in this section strongly support MERIT's choice of integrating Shapley values into the prototype~\cite{lundberg2017unified}
along with glass-box audit trails. We attribute the final score to CV, GitHub and LinkedIn coalitions, and also surface the 
full mathematical trail for each metric so that recruiters can audit each stage of the pipeline, rather than trust an explanation.

\section{Shapley Values for Data Source Attribution}
\label{sec:explainable-ai-shapley}

For MERIT to be considered as a "Glass-box" system, source-level attribution must be provided to provide a fair assessment to a recruiter as to how a candidate was scored.
Part of that required the system being able to capture the contribution of each data source to the final match score.
To achieve this, we focus on Shapley Values, a concept from Cooperative Game Theory that is used to attribute the winnings of a group of players to the contribution of each player.


\subsection*{Cooperative Game Theory in Recruitment}
To relate this concept to MERIT, we consider the final match score as "winnings", while the different data sources are considered the "players".
The main objective is to determine how each player contributed to the winnings.

The Shapley value for a data source $i$ is defined as the average marginal contribution of that source across all possible combinations (coalitions) of sources:
\begin{equation}
    \phi_i(v) = \sum_{S \subseteq N \setminus \{i\}} \frac{|S|! (n - |S| - 1)!}{n!} [v(S \cup \{i\}) - v(S)]
\end{equation}
Where:
\begin{itemize}
    \item $n$ is the total number of data sources.
    \item $S$ is a subset of sources that does not include source $i$.
    \item $v(S)$ is the score generated using only the sources in subset $S$.
\end{itemize}

\subsection*{Source Attribution and Fairness}
By calculating these values, MERIT can explicitly state the contribution of each data source to the final match score.
In other words, it is possible to determine whether a candidate got a high score due to self-reported claims on their CV, or from evidence-based metrics from a source, such as GitHub.
This level of transparency is critical for mitigating algorithmic bias and providing defensible hiring decisions.
By building on the SHAP (Shapley Additive Explanations) framework \cite{lundberg2017unified}, MERIT ensures that its source-level attribution remains mathematically consistent and locally accurate.

